\section{Availability and Licensing}
\label{sec:availability}

\paragraph{Code and artifacts.} Training, evaluation, auditing, Grad-CAM and
table-generation code are released with the paper, together with all evaluation
JSONs, the generated tables, and the figures. The frozen split manifest is
released as a CSV with the following SHA-256 digest:

\begin{center}\small\ttfamily
748322d62ace769445d0bd2363699097\\
0c7a5115a99c9344e139099c608c80de
\end{center}

\noindent No image, label file, or crop is redistributed; see
\texttt{DATA\_PROVENANCE.md} for the reason and for the exact source inventory.

\paragraph{Weights.} Three detector seeds are released, each with
\texttt{best.pt}, \texttt{last.pt} and an ONNX export, plus three stage-2
checkpoints and the original v1 run for reference, all with a SHA-256 manifest.
The seed-3/4 \texttt{yolov8s} and seed-0--2 \texttt{yolov8n} runs discussed in
Section~\ref{sec:results} were trained on the same frozen split under the same
schedule; their evaluation JSONs are committed so every number in this paper is
reproducible, but their weight files are not part of the released tree.
The detector checkpoints are derivative works of Ultralytics YOLOv8 and are
distributed under \textbf{AGPL-3.0}; the stage-2 checkpoints are fine-tuned
torchvision MobileNetV3-Large and are distributed under the same repository
licence for consistency. Third-party terms are stated in
\texttt{models/NOTICE.md}. The repository code and the audit artifacts are
AGPL-3.0.

\paragraph{Reproduction.} Regenerating every table requires no GPU and no
dataset: the committed JSONs are sufficient. The table script exits with an error
if a required field is absent rather than substituting a value, so a clean clone
either reproduces the tables in this paper or fails visibly.

\paragraph{Persistent identifier.} No DOI has been assigned at the time of this
draft. The manifest digest above is a content-level identifier; a Zenodo DOI and
exact release tag will be added to the camera-ready release.
